% ====================================
% intro-b.tex
% intro-b for rel-t-morse.tex
% historical remark
% ====================================

%Using microlocal sheaf theory, 
%one can investigate propagation phenemena of cohomology of sheaves. 

%In 1971, Zerner studied the analytic continuation of the solution of PDE in 
%\cite{Zer71}. 

Microlocal theory of sheaves lets us treat 
the propagation of cohomology of sheaves on manifolds 
in the language of cotangent bundles. 
In particular, microsupports of sheaves know 
the (co)direction-wise singularity. 

Note that given a sheaf (or a complex of sheaves) \(F\) 
on a real manifold of class \(C^\infty\) and a point \(p\) 
in its cotangent bundle, the point \(p\) does not belong 
to \(\muS(F)\), the microsupport of \(F\), 
if and only if for every degree \(j\), the \(j\)-th cohmology 
of \(F\) has no singularity along the (co)direction 
represented by \(p\). 

However, 
there are propagation phenomena which depend on 
the degree of cohomology. 
One cannot treat such phenomena by microsupports. 
\cite{KMS03a} introduced the notion of truncated microsupport. 
This has the information of (co)direction-wise 
singularity depending on the degree of cohomology. 
Both notions above have knowledge of microlocal nature. 
Then a question arises: can we collect this microlocal 
information and say something about 
the global propagation phenomena? 
One of its answers is given by the non-characteristic 
deformation lemma due to Kashiwara (Hol1). 
This has an origin from Zerner's so-called sweeping-out method 
we review later. 
What we give in this paper is a truncated version 
of the non-characteristic deformation lemma. 
As a result, we also give a truncated version of 
the microlocal Morse lemma for a morphism of manifolds. 
Finally we note that by considering the sheaf of hyperfunction 
solutions of a \(\mathcal{D}\)-module, 
one can obtain a global version of 
Holmgren's uniqueness theorem. 


\paragraph{Analytic point of view} 

In view of propagation, 
our results are along the lines of 
the Cauchy-Kowalevsky theorem. 
In 1971, Zerner showed a theorem of analytic extension of 
the solution of a single differential equation 
whose coefficients are holomorphic. 
In 1972, Bony-Schapira generalized the result to 
the case of families of differential operators. 
In 1975, Kashiwara used this method and 
the Mittag-Leffler argument for sheaves to extend the result 
to the case of solution complexes to a coherent \(\mathcal{D}\)-module. 
Note that Kashiwara used this method to show the constructibility theorem:  
if \(\mathcal{M}\) is a holonomic \(\mathcal{D}\)-module,
the solution complex \(\RHOM_{\mathcal{D}_{X}}(\mathcal{M},\mathcal{O}_{X})\) 
is constructible.

Our result in the case of hyperfunction solutions can be seen 
as a truncated version of these propagation.